\section{Discussion}

The results of this study provide compelling evidence for the "Invariance Breaking" hypothesis and the "Safe Haven" effect. By leveraging a longitudinal SCED approach, we were able to document, for the first time in a web-based sensing context, how the home environment fundamentally alters the human stress response.

\subsection{The "Safe Haven" as a Moderator}
The most striking finding is the disappearance of the Threat response at home. In the University lab, the "Threat" instruction worked as intended: it created urgency (faster RT) and freezing (reduced motion). The Bayes Factor of 66.7 leaves little doubt that the manipulation was effective in this context. 

However, at home, the exact same instruction—displayed on the exact same screen, in the same software—produced no such effect. Why?
From the perspective of Polyvagal Theory \cite{porges_safe_haven}, the home environment likely provides continuous, subconscious cues of safety—soft lighting, familiar smells, the absence of observers—that trigger "neuroception" of safety. This safety signal increases the threshold for sympathetic activation. Even when the cognitive system processes a semantic threat ("You are failing"), the autonomic nervous system, grounded in a safe context, applies a "vagal brake," dampening the fight-or-flight response. 

This has profound implications for Affective Computing. A stress-detection algorithm trained on lab data (where Stress = High Arousal) would fail in the home. It might detect "stress" where there is none (false positive), or conversely, fail to detect the subtle, regulated stress response characteristic of a safe environment (false negative).

\subsection{Context is a Lens, Not Noise}
In traditional engineering approaches, environmental variability is often treated as "noise" to be filtered out. Our results suggest the opposite: **context is a lens through which physiology must be interpreted.**

Consider a user with a heart rate of 80 bpm. 
\begin{itemize}
    \item In the \textbf{Lab (Evaluative Context)}, 80 bpm might represent a baseline state, rising to 100 bpm under stress.
    \item In the \textbf{Home (Safe Context)}, 80 bpm might represent a significant elevation from a relaxed baseline of 60 bpm.
\end{itemize}
A pervasive "universal threshold" (e.g., "HR > 90 = Stress") would misclassify both cases. Our SCED analysis shows that the \textit{reactivity slope} itself changes with context. The Home environment flattens the slope. Therefore, affective models must be parameterized by context: $Affect = f(Physiology, Context)$.

\subsection{Implications for Ubiquitous Computing Design}
Based on these findings, we propose three design principles for next-generation affective systems:

\begin{enumerate}
    \item **Calibrate per Context**: Systems should detect the user's semantic location (e.g., via Wi-Fi SSID or GPS) and switch between distinct "Context Profiles" (e.g., Focus Mode vs. Relax Mode) with different sensitivity thresholds.
    \item **Account for Dampening**: In safe environments, systems should expect lower signal-to-noise ratios in physiological reactivity. Algorithms may need to be more sensitive or rely on longer integration windows to detect affective shifts at home.
    \item **Prioritize Longitudinal Baselines**: One-shot calibration is insufficient. A user's baseline shifts daily and contextually. Continuous, unsupervised calibration—as mimicked by our longitudinal SCED approach—is necessary to capture the "User-in-Context."
\end{enumerate}

\subsection{Limitations}
While SCED provides high internal validity for the individual, the single-subject nature ($N=1$) limits population-level generalization. We cannot claim that \textit{all} users will show this specific Safe Haven effect; users with high trait anxiety might carry their "threat" everywhere. However, the goal of this study was to demonstrate the \textit{existence} of the phenomenon and validate the Camerala framework. Future work should deploy Camerala to a larger cohort ($N=50$) to quantify the inter-individual variance of this contextual modulation.
Additionally, while Motion and $E_{fluc}$ are robust, they are coarse proxies. Future iterations will integrate the rPPG signal more deeply once we improve its robustness to the specific dynamic lighting conditions found in the home.

\section{Conclusion}

We developed \textbf{Camerala}, a privacy-preserving web-based framework for longitudinal affective monitoring, and validated it with a 10-session SCED study. The study revealed a robust **Context $\times$ Appraisal Interaction**: the physiological and behavioral signatures of "Threat" observed in the laboratory were significantly dampened in the home environment.

These findings confirm the "Invariance Breaking" hypothesis: physiological signals do not have fixed meanings. They are dynamically constructed from the interaction of the body, the brain's appraisal, and the environmental context. As we move sensing from the lab to the living room, acknowledging this context-dependency is no longer optional—it is the prerequisite for valid affective computing. SCED provides the rigorous methodological framework needed to map these personalized, context-dependent landscapes.
